\documentclass[10pt,a4paper]{amsart}
\usepackage[margin=19mm]{geometry}
\usepackage{amsmath,amssymb,mathtools}
\usepackage[scr=rsfs,cal=euler]{mathalfa}
\usepackage{xfrac}
\usepackage[unicode,hidelinks,bookmarks=false]{hyperref}
\newcommand{\ie}{i.e.\ }
\newcommand{\cf}{cf.\ }
\newcommand{\resp}{respectively\ }
\newcommand{\Subsection}{\subsection}
\providecommand{\nobreakdash}{\nobreak\hbox{-}}
\setlength{\emergencystretch}{2em}
\multlinegap0pt
\usepackage{siunitx}
\usepackage{siunitx}
\usepackage{siunitx}
\usepackage{siunitx}
\usepackage{amssymb}
\usepackage{mathtools}
\usepackage{bm}
\usepackage{float}
\usepackage{siunitx}
\usepackage{cancel}
\usepackage[normalem]{ulem}
\newtheorem{proposition}{Proposition}
\title{On Higher-Order Geometric Refinements\\ of Classical Covariance Asymptotics:\\ an Approach via Intrinsic and\\ Extrinsic Information Geometry}
\author{Malik Amir, Sourangshu Ghosh}
\date{Source: arXiv:2604.12725v1 (CC BY 4.0)}
\begin{document}
\maketitle

\noindent\emph{Source and licence.} On Higher-Order Geometric Refinements\\ of Classical Covariance Asymptotics:\\ an Approach via Intrinsic and\\ Extrinsic Information Geometry. Version arXiv:2604.12725v1 (CC BY 4.0), distributed by arXiv under the Creative Commons Attribution 4.0 International licence, http://creativecommons.org/licenses/by/4.0/, as recorded in the arXiv licence metadata for this exact version and shown on the abstract page https://arxiv.org/abs/2604.12725v1. Full licence notice: \texttt{licenses/arxiv-2604.12725-CC-BY-4.0.txt}.\par\smallskip

\noindent\textit{Editorial scope.} Counted unit: the article's proposition prop:posterior\_MSE (source0.tex lines 1351--1360). The article's complete original proof of it is delivered with it (source0.tex lines 1362--1392, one \texttt{proof} environment of 31 lines). No mathematics is written, repaired or re-derived: the statement and the proof are the article's own lines, in the article's own notation, and no retained line is substituted or re-flowed.  Retained uncounted as the source-local context the proof needs: lines 1297--1310.  Nothing was omitted from any retained span, so the package carries no omission record.  No retained line contains a citation, so the wrapper carries no bibliography and no bibliography entry.  No retained line contains a figure environment, an image-inclusion command or drawing code, so no image is delivered and the package declares no asset dependency.  Editorial formatting only, in addition: an AMS \texttt{amsart} preamble that restates the article's own declarations and macros used by the retained lines, namely the environments \texttt{proposition} on separate counters, together with the standard packages those lines need.  The retained environments print in this document as Proposition 1 in order of appearance, whereas the article prints the same items with the numbering of its own full document.  The complete native source of the article as distributed in the arXiv e-print for this exact version is bundled unchanged as \texttt{sources/198422/source0.tex}.
\begin{proposition}[RLCT under the additive assumption]
\label{prop:RLCT_additive}
Under the additive normal crossing representation $K(\pi(u)) = \sum_{j=1}^r c_j u_j^{2k_j}$ with Jacobian $|\det D\pi(u)| = \prod_{j=1}^r |u_j|^{h_j} \cdot \varphi(u)$, the stochastic complexity integral
\[
Z_n := \int_{\Theta} \exp\big(-n K(\theta)\big) \, \varphi(\theta)\, d\theta,
\]
satisfies
\[
Z_n \sim C \, n^{-\lambda},
\qquad
\lambda = \sum_{j=1}^r \frac{h_j + 1}{2k_j},
\]
as $n \to \infty$, where $C > 0$ is a computable constant.
\end{proposition}

\begin{proposition}[Posterior MSE rate]
\label{prop:posterior_MSE}
Under the additive assumption with prior $\varphi(\theta) > 0$, the Bayesian posterior mean squared error satisfies
\[
\mathbb{E}_{\pi_n} \big[ \|\theta - \theta_0\|^2 \big]
\sim
C \, n^{-\min_j \frac{1}{k_j}},
\]
where $\pi_n(\theta) = \frac{\exp(-nK(\theta))\varphi(\theta)}{Z_n}$ is the posterior distribution.
\end{proposition}

\begin{proof}
Substituting $\theta = \pi(u)$ and using $\pi(u) - \theta_0 = \sum_{j=1}^r a_j u_j + O(\|u\|^2)$ for vectors $a_j \in \mathbb{R}^d$, we have
\[
\|\pi(u) - \theta_0\|^2 = \sum_{j=1}^r b_j u_j^2 + \sum_{j \neq \ell} b_{j\ell} u_j u_\ell + O(\|u\|^3),
\]
where $b_j = \|a_j\|^2 > 0$ and $b_{j\ell} = \langle a_j, a_\ell \rangle$. By symmetry of the integrand, the cross terms $u_j u_\ell$ (with $j \neq \ell$) vanish upon integration against the even function $\exp(-n \sum c_j u_j^{2k_j})$. Hence the numerator
\[
\int \|\theta - \theta_0\|^2 \exp(-nK(\theta))\varphi(\theta)\,d\theta
\sim
\sum_{j=1}^r b_j \underbrace{\int u_j^2 \exp(-nc_j u_j^{2k_j}) |u_j|^{h_j}\,du_j}_{=:I_j'(n)} \cdot \prod_{\ell \neq j} I_\ell(n).
\]
The scaling $t_j = n^{1/(2k_j)}u_j$ gives
\[
I_j'(n) = n^{-\frac{h_j+3}{2k_j}} B_j,
\qquad
B_j := \int t^2 \exp(-c_j t^{2k_j})|t|^{h_j}\,dt < \infty.
\]
Since $\frac{h_j+3}{2k_j} = \frac{h_j+1}{2k_j} + \frac{1}{k_j}$, the $j$-th term in the numerator scales as
\[
n^{-\frac{h_j+3}{2k_j}} \cdot \prod_{\ell \neq j} n^{-\frac{h_\ell+1}{2k_\ell}} = n^{-\lambda - \frac{1}{k_j}}.
\]
The dominant term corresponds to the smallest additional decay, i.e., $\min_j \frac{1}{k_j}$, giving
\[
\text{Numerator} \sim C_1 \, n^{-\lambda - \min_j \frac{1}{k_j}}.
\]
Dividing by $Z_n \sim C_0 \, n^{-\lambda}$ from Proposition~\ref{prop:RLCT_additive}, we obtain
\[
\mathbb{E}_{\pi_n}[\|\theta - \theta_0\|^2] \sim C \, n^{-\min_j \frac{1}{k_j}}.
\]
In the regular case $k_j = 1$ for all $j$, this gives $n^{-1}$. In the singular case with $k_j \geq 2$ for some $j$, the rate is $n^{-\min_j 1/k_j} \gg n^{-1}$.
\end{proof}

\end{document}
